\documentclass[10pt]{article}
\usepackage[margin=20mm]{geometry}
\usepackage{amsmath,amssymb,fontspec,unicode-math,xurl,hyperref}
\IfFontExistsTF{DejaVu Serif}{\setmainfont{DejaVu Serif}}{\setmainfont{Latin Modern Roman}}
\setmathfont{latinmodern-math.otf}
\setlength{\parindent}{0pt}\setlength{\parskip}{6pt}\sloppy
\begin{document}
Abstract: The proceedings paper, appendix, and pinned public code were inspected. Selecting recently recognized species does not guarantee that their visual concepts were absent from pretraining. This invalid inference affects the central benchmark interpretation, without establishing actual contamination. A separate discrepancy changes the published feature-filtering score: its unnormalized sum conflicts with the fixed threshold and released implementation. The headline Table 2 gains are arithmetically consistent. Scientific adjudication remains pending.\par
Author: AI-assisted technical review\par
Date:\par
Title: Technical Verification of Contrastive Visual Data Augmentation\par
\# Identity and scope\par
This report concerns Yu Zhou, Bingxuan Li, Mohan Tang, Xiaomeng Jin, Te-Lin Wu, Kuan-Hao Huang, Heng Ji, Kai-Wei Chang, and Nanyun Peng, Contrastive Visual Data Augmentation, ICML 2025, PMLR 267, pp. 78812–78826 [Retained latex: \textbackslash{}cite\{paper\}]. The identifier is 0181/I25-43550. The reviewed 15-page proceedings PDF includes its appendix. Page references are one-based PDF pages, agreeing with printed page labels.\par
The inspection covered NovelSpecies construction and interpretation, feature filtering, reported comparison arithmetic, Appendix A.2–A.4, and relevant code in the pinned release. The source was read statically; it was not executed. No model was trained or queried. Dataset images, pretraining corpora, and per-run predictions or human annotation records were not independently audited. Separate OpenReview attachment availability is unresolved in the source-preparation record.\par
\# Definitions and dependencies\par
CoDA extracts contrastive attributes, filters them, synthesizes and verifies images, and updates a target model. Table 2 tests SUN, iNaturalist, and NovelSpecies. The empirical contribution therefore extends beyond the novel-species interpretation.\par
For target and confusing classes \(C_T,C_M\), Section 3.2 defines a sum of pairwise CLIP-score ratios and uses a fixed cutoff of \(0.6\). Section 4 claims that selecting concepts discovered after the latest component-training cutoff guarantees novelty. Formal recognition date, first availability of relevant records, and actual training-corpus membership are distinct quantities.\par
\# Substantive findings\par
Problem 1 (Recognition date does not guarantee absence from pretraining).  The temporal implication underlying the asserted novelty guarantee is invalid. The supported impact is Major for this core benchmark interpretation; actual contamination remains unestablished.\par
\#\# Claim and location\par
The abstract and introduction (pp. 1–2), Section 4 (pp. 5–6), and Section 5.1 (p. 6) assert that species selected after a model’s cutoff were never encountered during training. Section 4 offers the date rule specifically as a guarantee despite unavailable corpora. Figure 1 and Section 3.1 motivate the approach with the clouded tiger cat, Leopardus pardinoides.\par
\#\# Complete check\par
Let \(t_r\) be recognition under the selected taxonomic status, \(t_p\) the first availability of relevant visual or descriptive records, and \(t_c\) the model’s latest training cutoff. The rule requires \(t_r>t_c\). It does not require \(t_p>t_c\). The ordering \[t_p<t_c<t_r\] satisfies the rule while allowing pretraining on an older record, including a photograph captioned with a historical subspecies name. Construct a model whose pre-cutoff training corpus contains that record. Its cutoff and the species’ recognition date meet the stated selection conditions, yet the underlying appearance has been encountered. Section 4 imposes no condition excluding this case.\par
The paper’s own example illustrates the distinction. The primary study published on 29 January 2024 groups previously recognized subspecies, including L. t. pardinoides and L. t. oncilla, as the clouded tiger cat, using historical museum and photographic material [Retained latex: \textbackslash{}cite\{taxonomy\}]. Recognition of a grouping is not creation of its underlying visual information. This evidence refutes the claimed implication; it does not identify an image in a particular closed training corpus.\par
\#\# Counterevidence, impact, and correction\par
A task of learning a newly adopted label or taxonomic distinction remains meaningful, as does low-resource adaptation. Those narrower interpretations support useful experiments but cannot provide the advertised unseen-concept guarantee. Conversely, unavailable corpora do not prove contamination.\par
NovelSpecies is an explicit principal contribution and the central unseen-concept test. The invalid implication thus affects a core interpretation. It does not invalidate the numerical accuracy tables, SUN/iNaturalist results, or the usefulness of CoDA. A supported correction is to describe recently recognized taxa under documented cutoffs and remove the unconditional guarantee. Establishing stricter novelty requires record dates, synonym-aware checks, and training provenance; no such comprehensive audit was performed here.\par
Problem 2 (The printed discriminability score differs from the released filter).  The published unnormalized sum fails its accompanying fixed-threshold interpretation. The supported finding is a local specification defect; experimental impact remains unresolved.\par
\#\# Claim and assumptions\par
Equation reference: eq:sum. Section 3.2, p. 4, defines \[
D=\sum_{i=1}^{n}\frac{a_i}{a_i+b_i},
\quad a_i=\operatorname{CLIP}(f,i_{C_T}),\quad
b_i=\operatorname{CLIP}(f,i_{C_M}),\] for equally many images per class. It states that scores below \(0.5\) favor the confusing class and filters at \(0.6\). Appendix A.3.2 identifies cosine similarity. The actual rendered page confirms the sum.\par
\#\# Calculation and source comparison\par
Take two images per class with \(a_i=b_i=0.2\). These positive cosine values have nonzero denominators and can be realized by unit vectors with first coordinate \(0.2\) against a unit text vector. Neither class is preferred, but ([eq:sum] (\#eq:sum)) gives \(D=1>0.6\). Five equal-score pairs give \(2.5\). With \(n=2,a_i=0.1,b_i=0.2\), every pair favors the confusing class, yet \(D=2/3>0.6\).\par
On the normal path where all image pairs are processed successfully, the pinned release [Retained latex: \textbackslash{}cite\{code\}], `filter.py`, lines 320–351, instead averages over all pairs after a softmax of scaled cosine scores: \[D_{\mathrm{release}}=\frac{1}{nm}\sum_{i=1}^{n}\sum_{j=1}^{m}
\frac{\exp(100a_i)}{\exp(100a_i)+\exp(100b_j)}.\] `generation.py`, lines 101–105, thresholds this statistic. Equal scores give \(0.5\) for every sample count, so that feature is rejected. Independent scalar calculations confirm the examples; no embeddings were computed and no author code was executed.\par
\#\# Counterevidence and impact\par
The probability-like interpretation and fixed thresholds suggest a normalized statistic, and the released code explicitly averages. That is substantial counterevidence against alleging that the experiments used the printed sum. Adding \(1/n\) alone also would not reproduce the code: softmax and all-pairs averaging differ from the published raw paired ratios. Raw cosine ratios need treatment of negative scores and zero denominators.\par
The finding is Minor as a demonstrated publication/release mismatch without an established failure of the operative experimental filter. The release already avoids the demonstrated sample-count failure. A precise documentation correction is the displayed release formula, if that was the experimental implementation. Historical commit provenance and benchmark effects remain Needs Evidence. This does not certify the entire pipeline.\par
\# Other checked claims and overall assessment\par
Table 2’s advertised 5:1 absolute gains check exactly: \[73.0-60.7=12.3,\quad83.4-78.3=5.1,\quad67.0-61.0=6.0\] percentage points. The setting uses one synthetic plus five real images per class. Table 3’s best LLaVA value, \(79.2\), is 18.0 points above \(61.2\). The relative calculation \(100(78.8-61.2)/61.2=28.758\%\) explains the roughly \(28.7\%\) improvement reported for a different comparison. Such arithmetic does not verify statistical significance; original run and annotation records remain unavailable.\par
The overall supported impact is Major for the unconditional novelty inference, with the score mismatch separately Minor. Actual contamination, historical training implementation, and effects of stricter novelty screening remain unresolved. The arXiv metadata identifies v2 (5 June 2025), whose abstract retains the guarantee. That was a metadata inspection, not a complete alternate-version comparison; no exhaustive erratum search was performed. Scientific adjudication remains pending; human reviews are empty. The issue grades below are unchanged. Adoption has not been requested or verified.\par
Issue assessment | Repair verification | Adoption | Remaining work\par
Issue assessment | Repair verification | Adoption | Remaining work\par
P1: Major; recognition date alone does not guarantee pretraining absence. | Logical scope correction only; actual contamination unestablished. | Not requested. | Document pre-cutoff visual/descriptive records and training coverage; audit novelty at the claimed scope.\par
P2: Minor; printed score and released filter differ. | Static formula and scalar diagnostics checked; no historical run repair. | Not requested. | Specify the experimental score and bind the historical commit/configuration.\par
9 Y. Zhou et al. Contrastive Visual Data Augmentation. PMLR 267 (2025), 78812–78826. https://proceedings.mlr.press/v267/zhou25h.html. Reviewed proceedings PDF. PlusLabNLP. CoDA, commit `a72e33a83fae0ee709663e7a9224d6beed9af2d9`, 16 June 2025. https://github.com/PlusLabNLP/CoDA. Static source inspection. T. G. de Oliveira et al. Ecological modeling, biogeography, and phenotypic analyses setting the tiger cats’ hyperdimensional niches reveal a new species. Scientific Reports 14 (2024), 2395. https://doi.org/10.1038/s41598-024-52379-8. Publisher HTML consulted. Y. Zhou et al. arXiv:2502.17709v2, 5 June 2025. https://arxiv.org/abs/2502.17709. Abstract and revision metadata only.\par
\# Evidence and roles\par
Original PDF SHA-256:\par
> .\par
The accompanying `verification/0181` contains actual source-hash checks, the release integrity ledger, independently written `independent\_checks.py` and its executed JSON output, a viewed rendering of source page 4, and preserved web-response evidence. The numerical script evaluates only scalar expressions and table entries.\par
The parent AI reviewer performed this review and integration. A separate native auxiliary challenge by rechecked the logical novelty claim, the taxonomy evidence already consulted by the parent, and the filter discrepancy. It confirmed the bounded findings: physical existence alone is not proof of an earlier visual record; a possible earlier visual or descriptive record defeats the universal guarantee, while actual training contamination remains unestablished. Its notes, source-page checks and hash bindings are retained. The parent’s final integration preserves that distinction. No human review, author acceptance, or empirical reproduction is asserted. Exact build status, report-PDF visual-QA status, paths, source/report hashes, and execution records are in `provenance.json`; unavailable operations are not marked completed.\par
\end{document}
